\section{Preliminaries}\label{sec:background}

\subsection{Notation and Polynomial Encoding}\label{subsec:prelim-notation}

We work over a finite field $\mathbb{F}$. All vectors, matrices, and tensors are
flattened in row-major order before commitment. For $n \in \mathbb{N}$, we write
$[n] := \{0,1,\ldots,n-1\}$.

All committed lengths are assumed to be powers of two. To make this precise, we
fix an integer $L \ge 1$ such that every committed object is smaller than $2^L$,
and is padded to a length
dividing $2^L$. We also fix a primitive $2^L$-th root of unity $\omega \in
\mathbb{F}$. For every power of two $n$ with $n \mid 2^L$, we define
\[
\omega_n := \omega^{2^L/n}
\qquad\text{and}\qquad
\mathbb{H}_n := \langle \omega_n \rangle.
\]
Thus $\mathbb{H}_n$ is the multiplicative subgroup of size $n$ used as the
evaluation domain for length-$n$ vectors. Its vanishing polynomial is
\[
\nu_n(X) := \prod_{a \in \mathbb{H}_n}(X-a) = X^n-1.
\]
Because all domains are derived from the same root $\omega$, we have
$\omega_{nm}^m=\omega_n$ whenever $n$, $m$, and $nm$ are powers of two dividing
$2^L$. This compatibility is used later when we align objects of different
shapes.

Let $\lambda_n^{(i)}(X)$ denote the $i$-th Lagrange basis polynomial on
$\mathbb{H}_n$.

For a vector $\mathbf{v}\in\mathbb{F}^n$, its polynomial encoding over
$\mathbb{H}_n$ is
\[
v(X) := \sum_{i \in [n]} \mathbf{v}[i] \lambda_n^{(i)}(X),
\]
where $\mathbf{v}[i]$ denotes the $i$-th entry of $\mathbf{v}$.

Throughout the paper, we need to manipulate the shapes
of a high-dimensional tensor. Given $\mathbf{V}\in \mathbb{F}^{a\times b\times c}$,
we write $\mathbf{V}_{\{abc\}}\in \mathbb{F}^{abc}$ for its flattened form,
and $\mathbf{V}_{\{a\times bc\}}\in \mathbb{F}^{a\times bc}$ for its flattened form with the first dimension preserved.
As a sanity check, note that
\[
\mathbf{V}_{\{a\times b\times c\}}[i, j, k] = \mathbf{V}_{\{a\times bc\}}[i, j \cdot c + k] = \mathbf{V}_{\{abc\}}[i \cdot bc + j \cdot c + k].
\]
A similar notation applies to tensors of other dimensions.
For tensors of the same shape, we write $\mathbf{V}+\mathbf{W}$ for their entrywise sum, and
$\mathbf{V}\odot\mathbf{W}$ for their Hadamard product.
We denote the size of tensor $\mathbf{V}$ by $|\mathbf{V}|$.

\begin{table}[t]
    \centering
    \small
    \caption{Main scalar parameters used throughout the paper.}
    \label{tab:main-scalars}
    \begin{tabular}{@{}p{0.20\columnwidth}p{0.72\columnwidth}@{}}
        \toprule
        Symbol & Meaning \\
        \midrule
        $\mathbb{F}$ & Finite field over which the proof is defined. \\
        $L$ & Global logarithmic domain bound; committed objects are padded to lengths dividing $2^L$. \\
        $T$ & Number of boosting rounds. \\
        $K$ & Number of classes; each round trains one tree per class. \\
        $M$ & Number of training samples. \\
        $d$ & Number of features. \\
        $B$ & Number of quantization bins; feature values lie in $[B]$. \\
        $h$ & Maximum root-to-leaf depth of the trees in the forest. \\
        $b$ & Number of public inference samples in one batch. \\
        $N_{\mathrm{tot}}$ & Total number of nodes in the encoded forest, including the dummy node. \\
        $N_{\mathrm{int}},N_{\mathrm{leaf}}$ & Number of internal nodes and leaf nodes in the forest. \\
        $\mathcal{B}$ & Public range bound used for fixed-point and comparison checks. \\
        \bottomrule
    \end{tabular}
\end{table}


\begin{table*}[t]
    \centering
    \small
    \caption{Main variables and tensor shapes used in the proof.}
    \label{tab:main-variables}
    \begin{tabular}{@{}p{0.22\textwidth}p{0.28\textwidth}p{0.42\textwidth}@{}}
        \toprule
        Variable & Shape & Meaning \\
        \midrule
        $\mathbf{D}$ & $[B]^{M\times d}$ & Quantized training dataset. \\
        $\mathbf{Y}$ & $\{0,1\}^{K\times M}$ & One-hot label matrix. \\
        $\mathbf{S},\mathbf{G}$ & $\mathbb{F}^{T\times K\times M}$ & Training scores and pseudo-residuals. \\
        $\mathbf{lc},\mathbf{rc}$ & $[N_{\mathrm{tot}}]^{N_{\mathrm{tot}}}$ & Left and right child indices of forest nodes. \\
        $\mathbf{N}^{\leftarrow},\mathbf{N}^{\rightarrow}$ & $[B+1]^{N_{\mathrm{tot}}\times d}$ & Lower and upper hyper-box bounds of nodes. \\
        $\mathbf{E}$ & $\{0,1\}^{N_{\mathrm{tot}}\times d}$ & One-hot split-feature mask of each node. \\
        $\boldsymbol{\Theta}$ & $\{0,1,\ldots,B\}^{N_{\mathrm{tot}}}$ & Split threshold index of each node. \\
        $\mathbf{tree},\mathbf{root},\mathbf{is\_leaf},\mathbf{score}$ & $[TK+1]^{N_{\mathrm{tot}}}$, $\{0,1\}^{N_{\mathrm{tot}}}$, $\{0,1\}^{N_{\mathrm{tot}}}$, $\mathbb{F}^{N_{\mathrm{tot}}}$ & Tree id, root flag, leaf flag, and leaf score. \\
        % $\mathbf{I}$ & $[B]^{b\times d}$ & Public inference input batch. \\
        % $\mathbf{U}^{\leftarrow},\mathbf{U}^{\rightarrow}$ & $\mathbb{F}^{TKb\times d}$ & Fetched leaf hyper-box bounds for inference. \\
        % $\mathbf{sc}$ & $\mathbb{F}^{TKb}$ & Fetched leaf scores for inference. \\
        % $\mathbf{Ind}$ & $[N_{\mathrm{tot}}]^{T\times K\times M}$ & Training-time leaf index reached by each sample. \\
        $\mathbf{C}_{\mathrm{leaf}},\mathbf{V}_{\mathrm{leaf}}$ & $\mathbb{F}^{N_{\mathrm{tot}}\times d\times B}$ & Leaf count and residual-sum histograms. \\
        $\mathbf{C},\mathbf{V}$ & $\mathbb{F}^{N_{\mathrm{tot}}\times d\times B}$ & Full count and residual-sum histograms after child aggregation. \\
        $\mathbf{C}_{\leq},\mathbf{V}_{\leq},\mathbf{C}_{>},\mathbf{V}_{>}$ & $\mathbb{F}^{N_{\mathrm{tot}}\times d\times B}$ & Prefix and suffix split statistics. \\
        % $\mathbf{A},\mathbf{A}'$ & $\mathbb{F}^{T\times K\times M}$ & Histogram addresses and random powers used in the histogram proof. \\
        % $\mathbf{P}$ & $\mathbb{F}^{N_{\mathrm{tot}}dB}$ & Random-power table for histogram aggregation. \\
        % $\boldsymbol{\Gamma},\boldsymbol{\chi}$ & $\mathbb{F}^{N_{\mathrm{tot}}\times d\times B}$, $\{0,1\}^{N_{\mathrm{tot}}\times d\times B}$ & Candidate split gains and validity mask. \\
        % $\boldsymbol{\Gamma}^{*},\mathbf{T}$ & $\mathbb{F}^{N_{\mathrm{tot}}}$, $\{0,1\}^{N_{\mathrm{tot}}\times B}$ & Selected gain and threshold mask for the chosen split. \\
        \bottomrule
    \end{tabular}
\end{table*}


\subsection{KZG Polynomial Commitments}\label{subsec:prelim-kzg}

We use KZG polynomial commitment scheme (PCS) for univariate polynomials~\cite{kzg10}.
The PCS consists of the following algorithms:

\begin{itemize}
    \item $\Setup(1^\lambda, d_{\max}) \rightarrow \pp$:
    Generates public parameters supporting commitments to polynomials of degree at most $d_{\max}$.

    \item $\Commit(\pp, f) \rightarrow c$:
    Commits to a polynomial $f$ and outputs a commitment $c$.

    \item $\Open(\pp, f, z) \rightarrow (v, \pi)$:
    Given a polynomial $f$ and an evaluation point $z$, outputs the evaluation
    $v = f(z)$ together with an opening proof $\pi$.

    \item $\Verify(\pp, c, z, v, \pi) \rightarrow \{0,1\}$:
    Verifies that $v$ is the correct evaluation of the polynomial committed in $c$
    at point $z$.
\end{itemize}

\paragraph{Batched openings.}
In addition to single-point openings, we use a batched opening interface that
allows the prover to open multiple evaluation claims together.

\begin{itemize}
    \item $\BatchOpen(\pp, \{f_i\}_{i \in [m]}, \{Q_i\}_{i \in [m]})
    \rightarrow (\{v_{i,z}\}_{i \in [m], z \in Q_i}, \Pi)$:
    For each committed polynomial $f_i$ and each query point set $Q_i$,
    outputs all claimed evaluations $v_{i,z} = f_i(z)$ together with a single
    batched proof $\Pi$.

    \item $\BatchVerify(\pp, \{c_i\}_{i \in [m]}, \{Q_i\}_{i \in [m]},
    \{v_{i,z}\}_{i \in [m], z \in Q_i}, \Pi) \rightarrow \{0,1\}$:
    Verifies that all claimed evaluations are consistent with the corresponding
    commitments.
\end{itemize}

\paragraph{Properties.}
We rely on the following standard properties. \emph{Completeness}: honestly generated openings
and batched openings always verify.  \emph{Binding}: after a
commitment is published, the prover cannot later open it to a different value at
the same point. 

\subsection{Gradient Boosted Decision Trees in Plaintext}\label{subsec:prelim-gbdt}

Given a commitment to a forest, the task of the proof of training is to show that
the forest arises from a valid execution of the GBDT training procedure.

Let $\mathbf{D}\in[B]^{M\times d}$ be the quantized training dataset, where
$M$ is the number of samples, $d$ is the number of features, and each feature
value lies in one of $B$ bins. Let $\mathbf{Y}\in\{0,1\}^{K\times M}$ be the
one-hot label matrix, with class-sample axes, for a $K$-class classification task. We treat the number
of boosting rounds $T$,
the number of classes $K$,
and the size of dataset $M$
as public parameters of the statement being proved.

\paragraph{Plaintext training.}
GBDT trains an additive model over $T$ rounds. In the proof-facing notation,
the cumulative scores and pseudo-residuals are arranged as tensors
$\mathbf{S},\mathbf{G}\in\mathbb{F}^{T\times K\times M}$, with axes ordered by
round, class, and sample. The entry $\mathbf{S}[t,k,r]$ denotes the cumulative
score of sample $r$ for class $k$ after round $t$ has been added, while
$\mathbf{G}[t,k,r]$ denotes the regression target used by the class-$k$ tree in
round $t$. For convenience, we write the virtual initial score state as
$\mathbf{S}[-1,k,r]=0$.

The class probabilities after round $t$ are obtained by applying softmax over
the class axis:
\begin{equation}
\boldsymbol{\pi}[t,k,r]
:=
\frac{\exp(\mathbf{S}[t,k,r])}
{\sum_{c\in[K]} \exp(\mathbf{S}[t,c,r])}.
\end{equation}
Let $\mathbf{Y}\in\{0,1\}^{K\times M}$ be the one-hot label matrix. The
multinomial logistic loss after round $t$ is
\begin{equation}
\mathcal{L}^{(t)}
:=
-\sum_{r\in[M]}\sum_{k\in[K]}
\mathbf{Y}[k,r]\log \boldsymbol{\pi}[t,k,r].
\end{equation}
Its partial derivative with respect to the score entry
$\mathbf{S}[t,k,r]$ is
\begin{equation}
\frac{\partial \mathcal{L}^{(t)}}{\partial \mathbf{S}[t,k,r]}
=
\boldsymbol{\pi}[t,k,r]-\mathbf{Y}[k,r].
\end{equation}
Therefore, after round $t$, the next negative-gradient target is
\begin{equation}\label{eq:plain-pseudoresidual}
\mathbf{G}[t+1,k,r]
\coloneqq
-\frac{\partial \mathcal{L}^{(t)}}{\partial \mathbf{S}[t,k,r]}
=
\mathbf{Y}[k,r]-\boldsymbol{\pi}[t,k,r],
\qquad t<T-1.
\end{equation}
The initial target is set as $\mathbf{G}[0,k,r]=\mathbf{Y}[k,r] - \mathbf{1}/K$. We call
$\mathbf{G}[t,\cdot,\cdot]$ the pseudo-residual slice for round $t$. As in gradient descent,
we want to move the score matrix toward the negative-gradient direction and thereby reduce the loss.
Thus in round $t$, the
learner trains one regression tree for each class $k\in[K]$ to fit the target
vector $\mathbf{G}[t,k,\cdot]\in\mathbb{F}^{M}$.\footnote{The above describes the training process for typical gradient-boosting decision trees. 
It is also possible to adapt our constructions to other variants,
e.g. XGBoost~\cite{xgboost16},
by modifying the training target and tree-splitting policy accordingly.
It is also possible to extend our approach to prove
the training and inference of random forests.
See \Cref{app:xgboost-extension,app:rf-extension}.}

\paragraph{Regression tree.}
Fix a round $t$ and a class $k$. The learner constructs a regression tree
$f^{(t,k)}:[B]^d\to\mathbb{F}$, which is a binary decision tree whose leaves
carry scalar prediction values. Each internal node $u$ is labeled by a feature
index $j_u\in[d]$ and a threshold bin $b_u\in[B]$. Given a sample
$\mathbf{x}\in[B]^d$, routing at node $u$ is determined by the predicate
$\mathbf{x}[j_u] < b_u$: if it holds, the sample is sent to the left child of
$u$; otherwise it is sent to the right child. The output
$f^{(t,k)}(\mathbf{x})$ is the value stored at the leaf reached by this routing
procedure.

Similarly, during training, each node $u$ is associated with the set
$I_u\subseteq[M]$ of samples currently reaching $u$. The root contains all
samples, i.e., $I_{\mathrm{root}}=[M]$, and each split partitions the parent
sample set into the left and right child sample sets.

\paragraph{How a node is split.}
For the class-$k$ tree in round $t$, let
\[
g_r := \mathbf{G}[t,k,r]
\]
denote the regression target attached to sample $r$. Consider a node $u$ with
current sample set $I_u$. For a candidate split specified by a feature
$j\in[d]$ and a threshold bin $b\in[B]$, define
\begin{align}
L_{u,j,b} &:= \{\, r\in I_u : \mathbf{D}[r,j] < b \,\},\\
R_{u,j,b} &:= \{\, r\in I_u : \mathbf{D}[r,j] \ge b \,\}.
\end{align}
The quality of this split is measured by the residual sum of squares (RSS):
\begin{equation}\label{eq:rss-plain}
\operatorname{RSS}(u,j,b)
=
\sum_{r\in L_{u,j,b}} (g_r-\bar{g}_L)^2
+
\sum_{r\in R_{u,j,b}} (g_r-\bar{g}_R)^2,
\end{equation}
where
\[
\bar{g}_L := \frac{1}{|L_{u,j,b}|}\sum_{r\in L_{u,j,b}} g_r,
\qquad
\bar{g}_R := \frac{1}{|R_{u,j,b}|}\sum_{r\in R_{u,j,b}} g_r
\]
are the empirical means on the two sides. Intuitively, the split is good if the
targets assigned to each child are well concentrated around a single mean.

Among all valid candidates, i.e. those for which both $L_{u,j,b}\neq\emptyset$ and $R_{u,j,b}\neq\emptyset$, the trainer selects a
feature-threshold pair minimizing \eqref{eq:rss-plain}.
If no valid split exists, the node is not split and becomes a leaf.

Once the tree structure is fixed, each leaf $\ell$ is assigned the mean target
value of the samples that reach it. Writing $I_\ell\subseteq[M]$ for the sample
set of leaf $\ell$, its output value is
\begin{equation}\label{eq:plain-leaf-value}
w_{\ell}^{(t,k)}
:=
\frac{1}{|I_\ell|}\sum_{r\in I_\ell} \mathbf{G}[t,k,r].
\end{equation}
Therefore, for every training sample $r$, the tree prediction
$f^{(t,k)}(\mathbf{D}[r,\cdot])$ is exactly the leaf value associated with the
leaf reached by sample $r$.\footnote{In our later implementation, the stored leaf score is the
learning-rate-scaled contribution $\mu\cdot w_{\ell}^{(t,k)}$ rather than the
unscaled mean $w_{\ell}^{(t,k)}$.}

\paragraph{Round update.}
After training all $K$ regression trees
$\{f^{(t,k)}\}_{k\in[K]}$ in round $t$, the learner updates the score matrix by
adding the tree outputs with learning rate $\mu>0$:
\begin{equation}\label{eq:plain-score-update}
\mathbf{S}[t,k,r]
=
\mathbf{S}[t-1,k,r] + \mu\, f^{(t,k)}(\mathbf{D}[r,\cdot]).
\end{equation}
This completes round $t$. The next round recomputes
$\boldsymbol{\pi}[t,\cdot,\cdot]$, forms the new pseudo-residual slice
$\mathbf{G}[t+1,\cdot,\cdot]$ using \eqref{eq:plain-pseudoresidual}, and trains a fresh
collection of $K$ regression trees against those updated targets. After $T$
rounds, the forest consists of the $TK$ trees trained all rounds.

At inference time, for a fresh sample $\mathbf{x}\in[B]^d$, the model
accumulates the per-class scores
\[
\sum_{t=0}^{T-1} \mu\cdot f^{(t,k)}(\mathbf{x})
\qquad\text{for each }k\in[K],
\]
and predicts from the resulting $K$-dimensional score vector, typically by
applying softmax.

\paragraph{Implication for the proof.}
The proof-critical step is the correctness of the regression-tree training at
each node, namely, that the committed split metadata and leaf values are
consistent with the pseudo-residual targets of that round and class. In
standard plaintext implementations, the optimizer often evaluates split
candidates by sorting samples according to each feature and scanning candidate
thresholds in that sorted order. In quantized systems, the same optimization is
commonly expressed through per-bin histograms. Our proof follows the latter
view: instead of proving a sorting procedure in zero knowledge, we prove the
correctness of the histogram statistics from which the trainer derives the same
split scores.

\subsection{Quantization}\label{subsec:prelim-quantization}

We follow the fixed-point quantization approach used in prior ZKP for ML works~\cite{sparrow24,zkdt20,zkboost26}. 
Given a real number $x$, we encode it as a signed integer
\[
q_x = \lfloor x \cdot 2^Q \rfloor,
\]
where $Q>0$ is a public quantization parameter. The corresponding fixed-point value is then interpreted as $x \approx q_x / 2^Q$.

Under this representation, addition is scale-preserving.
Multiplication, however, increases the scale and therefore requires an explicit
rescaling step. In particular, if
$x \approx q_x/2^Q$ and $y \approx q_y/2^Q$, then their product is represented by
\[
q_z = \left\lfloor \frac{q_x q_y}{2^Q} \right\rfloor,
\]
so that
\[
xy \approx \frac{q_z}{2^Q}.
\]

Throughout the paper, we assume the encoding of a fixed-point number
    is in range $[-\mathcal{B}/2, \mathcal{B}/2)$ for some public $\mathcal{B}>0$.
    We assume also that the field $\mathbb{F}$ is large enough to support all intermediate values without modular wraparound, so that the arithmetic over $\mathbb{F}$ corresponds to integer arithmetic on the fixed-point encodings.

\subsection{Table Lookup Argument}\label{subsec:prelim-lookup}
Given matrix $\mathbf{F}\in \mathbb{F}^{n\times m}$
    and matrix $\mathbf{T}\in \mathbb{F}^{N\times m}$,
    the lookup relation
    \[\mathbf{F}\subseteq \mathbf{T}\]
    means that every row of $\mathbf{F}$ appears as a row of $\mathbf{T}$.


In our construction, we need a variant of this argument,
    where the prover has a batch of $\mathbf{F}_i\in \mathbb{F}^{n\times m_i}$
    and $\mathbf{T}_i \in \mathbb{F}^{N\times m_i}$
    for $i\in [k]$,
    and wants to prove that
    \[\mathbf{F}_0\| \cdots \| \mathbf{F}_{k-1} \subseteq \mathbf{T}_0\| \cdots \| \mathbf{T}_{k-1}.\]
    This can be achieved by first padding each matrix to the same
    number of columns,
    and then randomly combining them with a verifier challenge,
    so it is reduced to a single-table lookup argument.
    \Cref{app:concatenated-table-lookup} gives the details of this argument and its properties.

Our construction mainly relies on the single-table lookup argument of \citet{cq++24},
    which reduces a table lookup argument to a basic lookup argument
    where the table is single-column.
    Additionally, in our construction, we use
    two types of (single) table lookup arguments:
    with preprocessing, and without preprocessing.
    
\paragraph{With preprocessing.}
Based on KZG PCS~\cite{kzg10}, \citet{cq++24} give a zero-knowledge argument for this relation
    where the online prover time is independent of $\mathbf{T}$.
    The key is to first reduce the problem to a single-column lookup, and then
    use logup~\cite{logup22} for proving the relation,
    while preprocessing $\mathbf{T}$ by storing the ``cached quotient''~\cite{cq22}.
    The drawback is that the preprocessing requires extra prover work,
    and the benefit is only significant when $\mathbf{T}$ is large.

\paragraph{Without preprocessing.}
For table lookup without preprocessing,
    we use the construction of \cite{cq++24} without running
    its preprocessing phase.
    The benefit for using this construction is that,
    as we will demonstrate later,
    we can batch all lookup instances for all (differently shaped) tables
    into one single lookup argument, and therefore achieve significant 
    reduction in proof size and prover time.
    The efficiency and proof size gain is enormous
    when the table is small.\footnote{
        We do not know whether it is possible to batch the lookup arguments with preprocessing, because the preprocessing is specific to each table and cannot be easily aligned across different tables.
        We leave this question for future work.
    }

